\documentclass[frenchb]{article}
\usepackage[T1]{fontenc}
\usepackage[latin1]{inputenc}
%Pour utilisation sous unix
%\usepackage[utf8]{inputenc}
%\usepackage[utf8x]{inputenc}
\usepackage{a4wide}
\usepackage{graphicx}
\usepackage{amssymb}
\usepackage{color}
\usepackage{babel}

\begin{document}

\begin{figure}[t]
\centering
\includegraphics[width=5cm]{inp_n7.png}
\end{figure}

\title{\vspace{4cm} \textbf{Rapport individuel - it�ration 3}}
\author{Jean-Baptiste Prevost\\ }
\date{\vspace{7cm} D�partement Sciences du Num�rique - Premi�re ann�e \\
2021-2022 }

\maketitle

\newpage
\tableofcontents

\newpage
\section{Travail effectu�}
Durant cette it�ration j'ai fait:\\
\begin{itemize}
\item Ajout de la sauvegarde des environnements.
\item Ajout choix de la carte pour lancer une partie.
\item Interface graphique jeu.
\item Cr�ation des classes du contr�leur pour les d�placements.\\
\end{itemize}

La sauvegarde et la s�lection des environnements est compl�tement fonctionnel. De plus, la modification des environnements sauvegard� est aussi disponible.
\\\\

\end{document} 